\begin{proof}
Write \(r:=|3|\) and \(\alpha:=-\log r>0\).  The element \(2\) is a
nonsquare in \(K\), and \(L/K\) is its unramified quadratic extension.  The
normalization of \(C\) is
\[
 \nu:\widetilde C:=\mathcal M(L\langle u\rangle)\longrightarrow C,
 \qquad y=u,\quad x=su.
 \tag{\(\heartsuit\)}
\]
It is finite, and it is an isomorphism over \(\{|y|>0\}\), since there
\(s=x/y\).  The remaining fibre is \(\mathcal M(L)\) over the unique
point \(o\) cut out by \(x=y=0\).  It is a singleton.  Thus \(\nu\) is
bijective on underlying Berkovich spaces; as \(K\) has characteristic zero,
this is exactly the normalization hypothesis in
\Cref{thm:p031-perfected-unibranch-curve-envelope}.

Put \(N:=\widetilde C\times_K B\), let \(z:=t/3\), and define
\[
 \widetilde f(u,t):=
 \max\left\{\log|s-z|,\ -\alpha\min\{|u|/r,1\}\right\}.
 \tag{\(\heartsuit_1\)}
\]
This is finite and continuous.  At \(u=0\), its second entry is zero and
\[
 \widetilde f(0,t)=\max\{\log|s-z|,0\}
                    =\log\max\{1,|z|\}.
 \tag{\(\heartsuit_2\)}
\]
The final expression is independent of a point of a normalization fibre:
if \(|z|\ne1\), the ultrametric inequality gives the equality immediately,
and if \(|z|=1\), both sides of \((\heartsuit_2)\) are zero even in the
cancellation case.  Away from \(u=0\), \(\nu\times\mathrm{id}_B\) is an
isomorphism.  Hence \(\widetilde f\) is constant on every fibre of the
finite surjection
\[
 \rho:=\nu\times\mathrm{id}_B:N\longrightarrow X.
\]
Finite surjective analytic maps are topological quotient maps, so
\(\widetilde f=f\circ\rho\) for one finite continuous function \(f\) on
\(X\).

For every \(m\ge1\), write uniquely
\[
 (s-z)^m=P_m(z)+sQ_m(z),\qquad P_m,Q_m\in K[z],
\]
and set
\[
 A_m:=yP_m(t/3)+xQ_m(t/3),
 \qquad
 h_m:=\frac1m\log\max\{|A_m|,|3|^m\}.
 \tag{\(\heartsuit_3\)}
\]
The \(A_m\) are global \(K\)-analytic functions on \(X\); the nonzero
constant \(3^m\) makes \(h_m\) finite.  Thus \(h_m\) is a global locally
Fubini--Study function, hence an allowed competitor type.  Pullback along
\(\rho\) gives
\[
 h_m\circ\rho=
 \max\left\{m^{-1}\log|u|+\log|s-z|,\ -\alpha\right\}.
 \tag{\(\heartsuit_4\)}
\]
Because \(|u|\le1\) on \(\widetilde C\), this is at most
\[
 \max\{\log|s-z|,-\alpha\}
 \leq\widetilde f.
\]
Hence \(h_m\leq f\).  The constant \(-\alpha\) also lies below \(f\), so
\(\mathcal A_f\) is nonempty.

Fix a root \(s_0\in L\) of \(s_0^2=2\).  The \(L\)-rigid point
\(b\in B\) with \(t=3s_0\) has completed residue field \(L\).  On the
closed unit \(L\)-disc, consider the two analytic maps to \(X\)
\[
 \gamma_\pm(u)=(\pm s_0u,u,3s_0).
\]
Their origins have the same image \(q=(o,b)\).  On the plus branch,
\[
 f\circ\gamma_+(u)=-\alpha\min\{|u|/r,1\}.
 \tag{\(\heartsuit_5\)}
\]
Let \(D_r=\{|u|\le r\}\) and let \(\eta_r\) be its Gauss point.  Every
\(g\in\mathcal A_f\) is classically psh, and so is its pullback to \(D_r\).
Since \(g(\gamma_+(\eta_r))\leq-\alpha\), the affinoid maximum principle
\cite[Theorem~6.4]{GR25Classical} gives
\[
              g\circ\gamma_+\leq-\alpha\quad\text{on }D_r.
\]
The constant competitor proves that the raw envelope equals \(-\alpha\) on
this disc, in particular at \(q\).

On the minus branch, \(|-s_0-s_0|=1\), and therefore
\[
 f\circ\gamma_-=0,
 \qquad
 h_m\circ\gamma_-=
 \max\{m^{-1}\log|u|,-\alpha\}.
 \tag{\(\heartsuit_6\)}
\]
For every \(u\ne0\), the latter tends to zero.  Thus the raw envelope is
zero at every point \(\gamma_-(u)\) with \(u\ne0\).  Its upper
semicontinuous regularization satisfies \(E\leq f\), because \(f\) is
upper semicontinuous.  Letting \(u=3^n\to0\) in the preceding equality now
gives
\[
 E(q)=0,
 \qquad E(\gamma_+(\eta_r))=-\alpha.
 \tag{\(\heartsuit_7\)}
\]
Indeed, the first equality follows from the approaching raw values and
\(f(q)=0\), and the second from the constant competitor and
\((\heartsuit_5)\).

If \(E\) were classically psh, then \(E\circ\gamma_+\) would be
subharmonic on \(D_r\).  Its value at the centre would be zero, whereas its
only Shilov-boundary value would be \(-\alpha\), contradicting the same
affinoid maximum principle.  Thus \(E\) is not classically psh.  Since every
CLD-v2 psh function is classically psh by \Cref{thm:cld-to-classical-psh}, it
is not CLD-v2 psh either.
\end{proof}
